\section{Time Evolution via Transverse Contraction}
\label{sec:transverse}

Section~\ref{sec:cft} established what we want to measure and why: the Loschmidt echo of a critical
quench carries conformal data, and that data sits in the spectrum of a transfer matrix that
propagates along the temporal direction. This section builds the object itself. Transverse
contraction has been developed in a series of works, beginning with the use of matrix product states
along the time direction for infinite chains~\cite{banuls2009, muller_hermes2012} and continued in
the studies of temporal entropies this thesis follows~\cite{carignano_marimon_tagliacozzo2024,
carignano_tagliacozzo2025}. Our calculations use the \texttt{ITransverse.jl}
implementation~\cite{carignano2025itransverse}, to which the reader interested in the software is
referred.

What follows is the bridge between the continuum statements of Section~\ref{sec:cft} and something a
computer can evaluate: the rotated network and the transfer matrix it defines, the temporal states
that play the role of boundary conditions, the transition matrix whose spectrum carries the
entropies, and the truncation that keeps the whole scheme consistent.

\subsection{The rotated network}
\label{sec:transverse:network}

Transverse contraction expresses the time evolution of a $D$-dimensional system as the contraction
of a $(D{+}1)$-dimensional tensor network, the extra dimension being time. For our chain this is a
two-dimensional network, a \ac{PEPS} up to Trotter error, whose tensors encode the entire dynamics.
Evaluating it is the difficulty: exact contraction of a \ac{PEPS} is \#P-hard in
general~\cite{schuch2007}, so any practical scheme is an approximation with a controlled bond
dimension.

The network is built from the ingredients of Equation~\eqref{eq:echo}. The initial and final states
are \acp{MPS}, and the evolution operator is a stack of \ac{MPO} layers, one per Trotter step. Since
the Hamiltonian is translation invariant and time independent, every bulk tensor $W$ of those layers
is the same object; Section~\ref{sec:mporep} constructs it for the \ac{NNN} model.

The usual approach contracts this network row by row, stacking Trotter layers onto the initial
\ac{MPS}, which is what makes the spatial entanglement grow until the simulation becomes
intractable. Instead we rotate the network by $90^\circ$ and contract it column by column. Under
this rotation the roles of the indices are exchanged: what were virtual spatial bonds become
physical temporal indices, and what were physical spatial indices become virtual temporal bonds. The
boundary columns are then \acp{tMPS}, $\bra{L}$ and $\ket{R}$, and a bulk column is a \ac{tMPO},
$\mathcal{E}$, acting as a transfer matrix in the spatial direction. By translation invariance the
same $\mathcal{E}$ appears at every spatial site.

This gives a concrete counterpart for each element of Section~\ref{sec:cft}. The conformal boundary
states at the two edges of the strip are the temporal \acp{MPS} $\bra{L}$ and $\ket{R}$; the strip
width is the number of temporal sites, the evolution time plus the cooling sites of the regulator
$\beta_0$; and the transfer matrix whose spectrum the \ac{CFT} predicts is the column $\mathcal{E}$.
The echo itself is recovered as the overlap $\braket{L|R}$, so contracting the network along space
is the same operation as evolving in time.

The two directions of the network have very different spectral character, which is worth making
explicit because the whole method depends on it. The rows are the Trotter layers
$U(\delta t)=e^{-iH\delta t}$, which are unitary, so their eigenvalues are pure phases. The column
$\mathcal{E}$ carries no such constraint: it is non-Hermitian, and its eigenvalues are generic
complex numbers. Prediction 4 of Section~\ref{sec:cft:temporal} concerns precisely this operator,
stating that after a critical quench its leading eigenvalues approach a common modulus, so that
$\mathcal{E}$ tends to a rescaled unitary at long times. The consequences of its non-Hermiticity for
the eigenvectors are developed at the end of Section~\ref{sec:mporep}, once the operator has been
constructed explicitly.

\subsection{Temporal states and the reduced transition matrix}
\label{sec:transverse:entropies}

Each application of $\mathcal{E}$ increases the bond dimension of the temporal states, so the
practical cost of the contraction is set by their temporal entanglement, the entanglement across a
cut in the time direction. Section~\ref{sec:cft:temporal} already gave the reason this is favourable
at a critical point: prediction 1 states that the generalized entropies grow only logarithmically
with the evolution time, so the required bond dimension grows polynomially rather than
exponentially. What that argument leaves open is which object one actually diagonalizes to obtain
those entropies, and that object is the subject of this subsection.

The answer is not the reduced density matrix of a single boundary. Because $\mathcal{E}$ is
non-Hermitian, its dominant left and right fixed points are different states, and the quantity we
want, the overlap $\braket{L|R}$, involves both. The relevant object is therefore the \ac{RTM},
built jointly from the two temporal boundaries,
\begin{equation}
    \mathcal{T}_t \propto \ket{R}\bra{L},
    \label{eq:rtm}
\end{equation}
normalized by $\braket{L|R}$ so that its eigenvalues sum to one. It is a transition matrix rather
than a density matrix: since $\bra{L}$ and $\ket{R}$ are not related by Hermitian conjugation,
$\mathcal{T}_t$ is non-Hermitian and its eigenvalues are complex. They are not probabilities, so the
Schmidt interpretation of an ordinary entanglement spectrum does not carry over.

The spectrum of $\mathcal{T}_t$ defines the family of generalized temporal entropies, the objects
introduced in high-energy physics and holography to describe transition matrices between distinct
states. Section~\ref{sec:cft} quoted the conformal prediction for these entropies; here we can give
their definition. The R\'enyi-2 member, which is the one used throughout this work, is
\begin{equation}
    S_{2}^{\mathrm{gen}}(t) = -\log \mathrm{Tr}\,(\mathcal{T}_t^2).
    \label{eq:s2gen}
\end{equation}
Its practical advantage is that $\mathrm{Tr}(\mathcal{T}_t^2)$ follows from contracting two copies
of the \ac{RTM}, with no diagonalization and no matrix logarithm. The von Neumann member
$S_1^{\mathrm{gen}}=-\mathrm{Tr}(\mathcal{T}_t\log\mathcal{T}_t)$ would require the full spectrum of
an exponentially large non-Hermitian matrix, which is why we avoid it here, as anticipated in
Section~\ref{sec:cft}. The same two-copy structure is what makes the R\'enyi-2 entropy the temporal
entropy that is accessible experimentally, through the two-replica protocol of
Ref.~\cite{boucomas2026measuring}.

In practice, once $\bra{L}$ and $\ket{R}$ have converged for a target time $T$, we evaluate
Equation~\eqref{eq:s2gen} at every internal bond of the temporal chain, each bond corresponding to a
cut $t\in(0,T)$. Since $\mathcal{T}_t$ is non-Hermitian the result is complex, and we plot its real
and imaginary parts separately against the normalized cut. One convention has to be fixed before
comparing with Equation~\eqref{eq:cft:chord}: our temporal cuts touch the edge of the strip at one
end, so the slope carries the boundary coefficient $\tfrac{c}{12}(1+1/n)$, equal to $c/8$ at $n=2$,
rather than the bulk value familiar from spatial entanglement. Every central charge quoted in
Section~\ref{sec:results} is read through this factor.

\subsection{Truncation: the transition matrix, not the density matrix}
\label{sec:transverse:truncation}

Since the bond dimension grows with every application of $\mathcal{E}$, the states must be truncated
at each step, and in the transverse picture there is a genuine choice of what to truncate on. The
conventional option discards the smallest singular values of the \ac{RDM} of a single boundary
state, which is the default in any tensor-network library. The alternative truncates on the \ac{RTM}
of Equation~\eqref{eq:rtm}, built from both boundaries at once. We use both, and the reason the
\ac{RTM} is our default comes down to three properties.

The first is gauge invariance. A tensor network can be re-gauged on any bond without changing the
physical overlap, and the two candidate objects respond differently: the isolated boundary \ac{RDM}
transforms by a congruence, which for the non-unitary transformations generic here distorts its
spectrum, whereas the \ac{RTM} transforms by a similarity, which leaves its eigenvalues untouched.
The Schmidt values of a single temporal boundary are therefore not physical, while the \ac{RTM}
spectrum is, and it is also the spectrum the \ac{CFT} predicts. The derivation is given in
Appendix~\ref{app:sub:rtm_truncation}.

The second is cost. The \ac{RTM} keeps only the part of each boundary that contributes to
$\braket{L|R}$, whereas the \ac{RDM} faithfully stores components of $\ket{R}$ orthogonal to
$\bra{L}$ that the final contraction discards. The bound on the \ac{RDM} cost is the state
entanglement, while for the \ac{RTM} it is the smaller operator-space
entanglement~\cite{carignano_marimon_tagliacozzo2024}, and in practice the difference in required
bond dimension is large enough to decide whether a long run fits in memory.

The third is stability, and here the advantage is reversed. Truncating a Hermitian positive operator
inherits the variational bound of standard tensor-network truncation, so the \ac{RDM} scheme
converges smoothly. The \ac{RTM} has no such guarantee, because for a non-Hermitian object singular
values and eigenvalues carry different information: components with large singular values can still
cancel in the trace that sets the physical overlap. Truncating the \ac{RTM} by singular values is
thus a well-motivated heuristic rather than a bound, and its convergence can be less smooth. The
argument, together with the local environment on which the truncation is actually performed, is
developed in Appendix~\ref{app:sub:rtm_truncation}.

The trade-off is therefore clear: the \ac{RTM} is gauge-invariant and much cheaper but not
guaranteed to converge monotonically, while the \ac{RDM} is rigorous and stable but expensive. We
take the \ac{RTM} as the production default, keep the \ac{RDM} as a fallback when the \ac{RTM}
output turns noisy, and verify in Section~\ref{sec:results} that the two agree wherever both
converge, so that no conclusion depends on the choice.

\subsection{The dynamics as an eigenproblem}
\label{sec:transverse:eigenproblem}

The rotation has one further consequence, and it is the one that determines how the calculation is
actually performed. Contracting the network along the spatial direction means applying the same
$\mathcal{E}$ over and over. For a long chain, the result is dominated by the leading eigenvectors
of that single operator: the contraction converges to the dominant left and right fixed points, and
the echo becomes their overlap. The dynamical problem has therefore been converted into an
eigenproblem for $\mathcal{E}$, which is what makes the whole approach worthwhile, since the cost no
longer grows with the length of the chain.

The natural way to solve it is the tensor-network version of the power method, absorbing one column
of the network into each boundary state at a time and truncating after every step. There is,
however, a difficulty specific to this problem. The power method isolates the dominant eigenvector
at a rate governed by the ratio $|\mu_1/\mu_0|$ of the two leading eigenvalues, so it works well
only when they are clearly separated. Prediction 4 of Section~\ref{sec:cft:temporal} says that a
quench to a critical point drives precisely this ratio towards one: the leading eigenvalues approach
a common modulus as $T$ grows. The regime in which the conformal physics we are after becomes
visible is thus exactly the regime in which a single-vector iteration stalls, mixing the dominant
eigenvector with its near-degenerate neighbours.

The remedy is to stop tracking one vector and track the whole cluster instead, iterating a block of
$k$ states that spans the degenerate subspace. Convergence is then set by the separation between the
block and the rest of the spectrum, which stays finite even when the eigenvalues inside the block do
not separate. Because $\mathcal{E}$ is non-Hermitian, every ingredient of such a method, from
orthogonalization to the notion of overlap itself, has to be rebuilt in a biorthogonal language.
This block power method is the main algorithmic development of this work, and
Section~\ref{sec:numerics} presents it in full.
